\documentclass[11pt]{article}
\usepackage[margin=1in]{geometry}
\usepackage{mathpazo}
\usepackage{amsmath,amssymb}
\usepackage{microtype}
\usepackage{booktabs}
\usepackage{graphicx}
\usepackage[dvipsnames]{xcolor}
\usepackage{tikz}
\usetikzlibrary{positioning,arrows.meta,calc,fit,backgrounds,decorations.pathreplacing}
\usepackage[most]{tcolorbox}
\usepackage{enumitem}
\usepackage{listings}
\usepackage[font=small,labelfont=bf]{caption}
\usepackage{hyperref}
\hypersetup{colorlinks=true,linkcolor=MidnightBlue,citecolor=MidnightBlue,urlcolor=MidnightBlue}
\usepackage{cleveref}
\urlstyle{tt}
\def\UrlBreaks{\do\_\do/\do.}
\newcommand{\file}[1]{\nolinkurl{#1}}
\newcommand{\code}[1]{\texttt{#1}}
\newcommand{\mhh}{m_{hh}}
\newcommand{\bC}{\boldsymbol{C}}
\newcommand{\snoise}{s_{\text{noise}}}

\definecolor{trainable}{HTML}{1B6CA8}
\definecolor{physics}{HTML}{2E933C}
\definecolor{datac}{HTML}{7B4EA3}
\definecolor{warn}{HTML}{C0392B}

\newtcolorbox{keyidea}[1][]{enhanced,breakable,colback=physics!6,colframe=physics!70!black,
  fonttitle=\bfseries,title={#1},boxrule=0.6pt,arc=2pt,left=6pt,right=6pt}
\newtcolorbox{aside}[1][]{enhanced,breakable,colback=trainable!5,colframe=trainable!70!black,
  fonttitle=\bfseries,title={#1},boxrule=0.6pt,arc=2pt,left=6pt,right=6pt}
\newtcolorbox{assume}[1][]{enhanced,breakable,colback=datac!5,colframe=datac!70!black,
  fonttitle=\bfseries,title={#1},boxrule=0.6pt,arc=2pt,left=6pt,right=6pt}
\newtcolorbox{watch}[1][]{enhanced,breakable,colback=warn!5,colframe=warn!70!black,
  fonttitle=\bfseries,title={#1},boxrule=0.6pt,arc=2pt,left=6pt,right=6pt}
\setlist{itemsep=2pt,topsep=3pt}
\setlength{\emergencystretch}{3em}

\lstset{language=Python, basicstyle=\ttfamily\small, columns=fullflexible,
  keepspaces=true, commentstyle=\color{physics!70!black}\itshape,
  keywordstyle=\color{trainable!90!black}\bfseries, stringstyle=\color{datac},
  showstringspaces=false, frame=single, rulecolor=\color{gray!40},
  backgroundcolor=\color{gray!3}, xleftmargin=4pt, xrightmargin=4pt,
  aboveskip=6pt, belowskip=6pt}

\title{\textbf{The paper appendices}\\[4pt]
\large Plan: Appendix A, the smooth model; Appendix B, the neural-network density model}
\author{P.~Bhat (Fermilab), N.~De~Filippis (Politecnico di Bari),\\
K.F.~Johnson, H.B.~Prosper, R.~Starnes (Florida State University)}
\date{6 November 2026}

\begin{document}
\maketitle

\begin{abstract}
\noindent
The current draft appendix (\file{docs/paper/appendix_uncertainty.tex}, v4) describes the smooth model
but still carries the old networks, ensembles and the Bayesian network as cross-checks, uses the word
``discrepancy'', and has figures with the old networks in them. We propose two self-contained
appendices for the paper, describing only the two models to be published: \textbf{Appendix A}, the
smooth spline model, which is the model of the paper; and \textbf{Appendix B}, the neural-network
density model, as an independent model of the same quantity. Everything else (the bin-value networks,
ensembles, the Bayesian network, the Laplace mixture) is left out. The plan fits a two-week schedule;
decisions for you are in \cref{sec:decisions}.
\end{abstract}

\tableofcontents

% ===========================================================================
\section{Principles}

\begin{itemize}
  \item \textbf{Two models only.} Each appendix describes one model completely: what it is, how it is
    fitted, its uncertainty, and how well it does. B refers to A for everything they share (the data,
    the measures, the model-error term) and does not repeat it.
  \item \textbf{One language.} The likelihood-ratio statistic $-2\ln\lambda$ for accuracy; pulls for
    calibration; ``model-error term'' for the systematic uncertainty, introduced in A as a systematic
    uncertainty in the physicists' sense. ``Discrepancy'' and ``deviance'' disappear.
  \item \textbf{Every number from a script in the repository.} Each figure and table is regenerated
    from the committed model files (\file{runs/SMOOTH_J40}, \file{runs/FTD_SiLU_huber_d2_pearson_curv1e-06}),
    with no number copied from an old workspace.
  \item \textbf{For the reader of the paper}: a particle physicist who knows HEFT and Poisson statistics
    but not splines or quadrature. Concise, with the pedagogy of the notes condensed into one sentence
    each where needed.
\end{itemize}

% ===========================================================================
\section{Appendix A: the smooth model}
\label{sec:A}

Target length: 5--6 journal pages. Sources: \file{dihiggs_model.pdf} (sections 3--8), the current
appendix, \file{smooth_results.pdf}.

\begin{enumerate}[label=\textbf{A.\arabic*}]
  \item \textbf{The form and the data.} The 23-term form $d\sigma/dm = \sum_i C_i(\theta)a_i(m)$ and when it
    is exact (no $\theta$ dependence anywhere else in the calculation). The POWHEG spectra as exact Poisson
    counts ($t_{sk} = q_s n_{sk}$); the fit rows, the quick-test rows and the 1,000 test spectra.
  \item \textbf{The model.} Cubic B-splines times the threshold factor, knots uniform in $\sqrt{m-250}$;
    the prediction is the bin integral, linear in the 920 coefficients. The free fit as the limit with
    one box per bin, which sets the floor. \emph{Figure: \file{fig_basis}.}
  \item \textbf{Fit and posterior.} Poisson maximum likelihood by iterated weights from the model (why not
    the observed counts, in one sentence); the second-difference prior with strengths from the evidence;
    the exact Gaussian posterior; its cost.
  \item \textbf{The systematic uncertainty.} The three uncertainties (Poisson noise, fit, model error); the
    floor reached by both the smooth model and the free fit; its growth with cancellation and its smoothness
    in $m$; why the errors must vary from run to run; the form $\varepsilon(m)\sum_i|C_iA_{ik}|$ as the
    largest spread relative errors $\varepsilon$ can produce; $\varepsilon$ from the quick-test rows (3
    parameters); what it does not include (scales, PDFs, top-mass scheme). \emph{Figure:
    \file{fig_model_error} (from \file{dihiggs_model.pdf}): $\varepsilon(m)$, pull width against $m$,
    $\langle z^2\rangle$ by cancellation fifth for four treatments.}
  \item \textbf{Validation on the test spectra.} Accuracy (table: smooth model and free fit;
    figure against $m$), calibration (table), three example spectra (\file{fig_spectra_dz}, replacing the
    nine-panel figure), inside a bin (merged-bin test), the edges (240--255 and 1455--1500~GeV).
  \item \textbf{The coefficient functions.} The 23 $a_i(m)$ with their exact bands, \emph{without} the old
    network's points (new figure).
  \item \textbf{Using the model.} What is published (the model file, the code), the valid range
    250--1455~GeV, the uncertainty for one bin and for sums of bins, and the SM example
    ($28.11$~fb in 255--1455~GeV; fit uncertainty 0.006~fb; model-error term 0.03--0.17~fb).
\end{enumerate}

% ===========================================================================
\section{Appendix B: the neural-network density model}
\label{sec:B}

Target length: 3--4 journal pages. Sources: \file{density_plan.pdf} (with addendum),
\file{density_results.pdf}.

\begin{enumerate}[label=\textbf{B.\arabic*}]
  \item \textbf{Why a second model.} An independent representation of the same 23 functions, built in an
    entirely different way, tests whether the results of A depend on the choice of splines and prior.
  \item \textbf{The network.} $a_i(m) = \phi(m)\,P_i(z)\,e^{Q_i(z)}$: $P$ a six-layer network with SiLU
    activations and a $\tanh$ output, $Q$ linear, input $z = 2\sqrt{(m-250)/1205}-1$; 3,944 parameters.
  \item \textbf{Bin integrals by Gauss--Legendre quadrature.} Eight nodes per bin; why the quadrature must be
    far more accurate than the noise (its error becomes a bias in the density); the cost (640 evaluations per
    step, independent of the batch). One paragraph on why a network trained on bin contents at bin centres
    is not a density (biased by about 1\% near threshold on a realistic toy), without describing the old
    networks.
  \item \textbf{Training.} L1 loss, then the Huber loss of the pulls with Poisson variances from the model;
    the smoothness penalty (curvature of $P_ie^{Q_i}$ in $z$, $\lambda = 10^{-6}$) and why it is needed (the
    poorly determined $a_{16}$--$a_{19}$).
  \item \textbf{Results.} Accuracy (table: A, free fit, B); calibration with its own $\varepsilon$, which
    equals A's: the systematic is a property of POWHEG, not of either model; agreement of the densities
    (0.31\% median, equal at and between bin centres); the coefficient functions against A's bands.
    \emph{Figures: accuracy against $m$ for A and B; $\varepsilon(m)$ of both;
    $d\sigma/dm$ of A and B for three spectra with the ratio B/A below (new, without the old networks); the
    23 $a_i$ of B over A's bands.}
  \item \textbf{Limitations.} A single network (no fit uncertainty of its own, unless the seeds ensemble is
    done); slightly less accurate than A just above threshold.
\end{enumerate}

% ===========================================================================
\section{Files}

\begin{itemize}
  \item \file{docs/paper/appendix_A_smooth.tex}, \file{docs/paper/appendix_B_density.tex}: the two appendices,
    each to be \code{\textbackslash input} after \code{\textbackslash appendix} in the paper (JHEP style),
    with \file{appendix_refs.bib}.
  \item \file{docs/paper/appendices_draft.tex}: a stand-alone wrapper that compiles both, for reading.
  \item Figure scripts in \file{docs/paper/}, one per figure, all reading the repository's model files through
    \file{smooth_common.py}. New: \file{fig_A_coefficients.py}, \file{fig_A_merged.py} (ported from
    \file{docs/smooth/scripts/merged.py}), \file{fig_B_accuracy.py}, \file{fig_B_density.py},
    \file{fig_B_coefficients.py}, \file{fig_B_epsilon.py}. Reused: \file{fig_basis.py},
    \file{fig_spectra_dz.py}, \file{docs/model/model_error_study.py}.
  \item The current appendix (v4) moves to \file{docs/paper/v4_smooth/}, like v2 and v3; the executive summary
    is updated at the end.
\end{itemize}

% ===========================================================================
\section{Schedule (two weeks)}

\begin{center}
\small
\begin{tabular}{ll}
\toprule
days 1--3 & Appendix A: text, new figures, the merged-bin script ported; you review \\
days 4--5 & Appendix B: text and figures; you review \\
in parallel (5090) & optional: seeds ensemble and merged-bin test for B (\cref{sec:decisions}) \\
days 6--8 & revisions from you and the colleagues; references; consistency with the main text \\
days 9--10 & executive summary; final numbers; tagged commit \\
\bottomrule
\end{tabular}
\end{center}

% ===========================================================================
\section{Decisions for you}
\label{sec:decisions}

\begin{keyidea}[Decisions before I proceed]
\begin{enumerate}
  \item \textbf{Structure}: A and B as in \cref{sec:A,sec:B}. Agreed?
  \item \textbf{The old networks}: left out entirely, except one sentence in B.3 that a network trained on bin
    contents at bin centres does not give the density (with the toy result). Agreed, or leave even that out?
  \item \textbf{Seeds ensemble for B} (10 seeds of L1 + fine-tune on the 5090, a few hours): gives B a fit
    uncertainty and a measured spread between bin centres. Worth it for the paper, or state the limitation?
  \item \textbf{Merged-bin test for B} (one L1 + fine-tune on 30~GeV bins): makes B's claim to be a density as
    strong as A's. Worth it?
  \item \textbf{The main text}: who writes it, and which section labels should the appendices refer to (now
    \code{sec:model}, a placeholder)? Should the appendices assume the main text has already introduced the
    23-term form and the POWHEG samples, or repeat them briefly?
  \item \textbf{Document class}: JHEP (\code{jheppub}) for the wrapper too, or a plain article for drafts?
\end{enumerate}
\end{keyidea}

\end{document}
